% Vietnamese translation for OI Public Library
% Source-File: IOI中国国家候选队论文/2005/唐文斌/唐文斌.doc
\documentclass[11pt]{article}
\usepackage{oipl-vn}

\newcommand{\Oh}{\mathcal{O}}
\newcommand{\st}{\mathrel{\mathrm{s.t.}}}

\title{Khó theo chiều thuận thì xét chiều ngược: bàn về ứng dụng của tư duy nghịch hướng trong giải bài}
\author{Tang Wenbin}
\date{2005}

\begin{document}
\maketitle

\origtitle{Khó theo chiều thuận thì xét chiều ngược: bàn về ứng dụng của tư duy nghịch hướng trong giải bài}
\sourcepath{IOI China candidate team papers / 2005 / Tang Wenbin.doc}

\renewcommand{\abstractname}{Tóm tắt}
\begin{abstract}
Tư duy nghịch hướng là một cách suy nghĩ về bài toán. Nó đi ngược thói quen
thông thường, và chính đặc điểm ấy khiến nhiều bài toán không thể giải, hoặc
rất khó giải, bằng lối suy nghĩ thuận trở nên sáng rõ. Bài viết dùng một số
ví dụ để tổng kết các ứng dụng của tư duy nghịch hướng trong giải bài tin học.
\end{abstract}

\noindent\textbf{Từ khóa:} tư duy nghịch hướng; nguyên lý bù trừ; tìm kiếm
tham số; nhị phân; quy hoạch động; ghi nhớ hóa.

\section{Mở đầu}

Trước hết xét một bài toán đơn giản.

\begin{quote}
Trên mặt phẳng có bốn điểm tạo thành một hình vuông cạnh $1$. Ta được thực
hiện thao tác sau: mỗi lần chọn hai điểm $A$ và $B$, lấy điểm $C$ đối xứng
với $A$ qua $B$, rồi bỏ $A$ đi.

Hãy chứng minh rằng sau hữu hạn thao tác không thể thu được một hình vuông có
cạnh lớn hơn $1$.
\end{quote}

Kết quả sau các thao tác khá phức tạp. Nếu đi theo hướng thuận, việc chứng
minh mệnh đề rất khó. Ta thử nhìn từ chiều ngược lại: dễ thấy mỗi bước thao
tác đều khả nghịch. Nói cách khác, nếu ta có thể làm hình vuông lớn lên, thì
cũng có thể làm nó nhỏ đi.

\paragraph{Chứng minh.}
Không mất tính tổng quát, giả sử bốn đỉnh ban đầu đều là điểm nguyên. Giả sử
có thể qua hữu hạn thao tác thu được một hình vuông cạnh lớn hơn $1$. Khi đảo
ngược toàn bộ các thao tác ấy và áp dụng lên hình vuông ban đầu, ta sẽ thu
được một hình vuông cạnh nhỏ hơn $1$.

Vì bốn đỉnh đều là điểm nguyên, sau mỗi thao tác tọa độ các điểm vẫn là số
nguyên. Do đó hình vuông cuối cùng vừa có cạnh nhỏ hơn $1$, vừa có bốn đỉnh
đều là điểm nguyên. Điều này hiển nhiên không thể xảy ra. Vậy giả thiết sai,
mệnh đề được chứng minh.

Ví dụ trên cho thấy ứng dụng của tư duy nghịch hướng trong toán học. Khi con
đường trước mặt tưởng như bế tắc, đổi góc nhìn có thể làm bài toán mở ra ngay.

\section{Ví dụ 1: Dinner Is Ready}

\subsection{Tóm tắt đề bài}

Mẹ nấu $M$ khúc xương để chia cho $n$ đứa trẻ. Đứa trẻ thứ $i$ có hai tham số
$\mathit{Min}_i$ và $\mathit{Max}_i$, lần lượt là số xương ít nhất và nhiều
nhất mà em ấy muốn nhận.

Dữ liệu vào gồm $n$ và $M$ với $0<n\le 8$, $M>0$, sau đó là $n$ dòng chứa
$\mathit{Min}_i,\mathit{Max}_i$ với
$0\le \mathit{Min}_i\le \mathit{Max}_i\le M$. Cần in ra số cách phân phối sao
cho không lãng phí khúc xương nào.

\subsection{Phân tích ban đầu}

Mô hình của bài toán rất đơn giản: đếm số nghiệm nguyên của hệ
\[
\begin{cases}
  \displaystyle\sum_{i=1}^{n} X_i=M,\\
  \mathit{Min}_i\le X_i\le \mathit{Max}_i,\qquad 1\le i\le n.
\end{cases}
\]
Ta biết hệ đơn giản
\[
  \sum_{i=1}^{n} X_i=M,\qquad X_i\ge 0
\]
có số nghiệm nguyên là
\[
  \binom{M+n-1}{n-1}.
\]

Đặt $Y_i=X_i-\mathit{Min}_i$. Khi đó hệ ban đầu trở thành
\[
\begin{cases}
  \displaystyle\sum_{i=1}^{n} Y_i
  = M-\displaystyle\sum_{i=1}^{n}\mathit{Min}_i,\\
  0\le Y_i\le \mathit{Max}_i-\mathit{Min}_i,\qquad 1\le i\le n.
\end{cases}
\]
Giới hạn dưới đã được xử lý bằng phép đổi biến, nhưng giới hạn trên vẫn khiến
ta chưa thể tính đáp án trực tiếp.

\subsection{Dùng tư duy nghịch hướng}

Gọi $S$ là tập toàn cục, tức tập các nghiệm nguyên thỏa $X_i\ge \mathit{Min}_i$.
Gọi $S_i$ là tập các nghiệm trong $S$ thỏa ràng buộc
$X_i\le \mathit{Max}_i$, và gọi $\overline{S_i}$ là phần bù của $S_i$ trong
$S$, tức các nghiệm thỏa $X_i>\mathit{Max}_i$.

Đích cần tính là $S_1\cap S_2\cap\cdots\cap S_n$, nhưng tập này khó đếm trực
tiếp. Trái lại, các tập bù như $\overline{S_i}$ và giao của chúng lại có thể
đếm được bằng cùng công thức nghiệm không âm sau khi trừ đi phần vượt quá
giới hạn trên. Vì vậy ta chuyển phép tính mục tiêu sang các tập bù:
\[
\begin{aligned}
\mathit{Answer}
&= |S_1\cap S_2\cap\cdots\cap S_n|\\
&= |S|
 - \sum_i |\overline{S_i}|
 + \sum_{i<j} |\overline{S_i}\cap\overline{S_j}|
 - \cdots
 + (-1)^n |\overline{S_1}\cap\overline{S_2}\cap\cdots\cap\overline{S_n}|.
\end{aligned}
\]
Đây chính là một dạng của nguyên lý bù trừ. Đến đây bài toán đã được giải:
ta không đếm trực tiếp tập đáp án khó xử lý, mà đếm các tập bù dễ xử lý hơn.
Thuật toán dựa trên bù trừ có độ phức tạp
\[
  \Oh(2^n(n+M)).
\]

\section{Ví dụ 2: Greedy Path}

\subsection{Tóm tắt đề bài}

Có $n$ thành phố và $m$ tuyến đường có hướng giữa chúng. Khi một du khách đi
từ thành phố $i$ đến thành phố $j$, chi phí phải trả là $C_{i,j}$ và thời
gian đi là $T_{i,j}$. Công ty du lịch chỉ có một xe, nên giám đốc muốn tổ
chức một hành trình vòng để thu được nhiều tiền nhất.

Với một chu trình $G$, hàm tham lam $F(G)$ bằng tổng chi phí chia cho tổng
thời gian. Hãy tìm một chu trình làm $F(G)$ lớn nhất. Nếu không có chu trình
thì in $0$; nếu có nhiều đáp án, in tùy ý một chu trình tối ưu.

\subsection{Phân tích ban đầu}

Bài toán yêu cầu tìm một chu trình, nhưng mục tiêu không phải là tổng trọng
số lớn nhất hoặc nhỏ nhất, nên không thể áp dụng trực tiếp các thuật toán
kinh điển.

Gọi đồ thị là $G=(V,E)$, và gọi $S$ là tập mọi chu trình
$C=(V',E')$ của $G$. Ta cần tìm một chu trình trong $S$ làm
\[
  F(C)=
  \frac{\displaystyle\sum_{e\in E'} C_e}
       {\displaystyle\sum_{e\in E'} T_e}
\]
đạt giá trị lớn nhất.

\subsection{Dùng tư duy nghịch hướng}

Nếu đã biết $C^*=(V^*,E^*)\in S$ là chu trình tối ưu, thì
\[
  t^* = F(C^*) =
  \frac{\displaystyle\sum_{e\in E^*} C_e}
       {\displaystyle\sum_{e\in E^*} T_e}.
\]
Từ đó ta định nghĩa
\[
  o(t)=\max_{C=(V',E')\in S}
  \left(\sum_{e\in E'} C_e - t\sum_{e\in E'} T_e\right).
\]

Ta đưa ra phỏng đoán: nếu $o(t^*)=0$, và tồn tại
$C^*=(V^*,E^*)\in S$ sao cho
\[
  \sum_{e\in E^*} C_e - t^*\sum_{e\in E^*} T_e=0,
\]
thì $C^*$ là một chu trình tối ưu.

\paragraph{Chứng minh.}
Giả sử tồn tại một chu trình khác $C_1=(V_1,E_1)\in S$ tốt hơn $C^*$. Khi đó
\[
  \frac{\displaystyle\sum_{e\in E_1} C_e}
       {\displaystyle\sum_{e\in E_1} T_e}>t^*,
\]
suy ra
\[
  0<
  \sum_{e\in E_1} C_e - t^*\sum_{e\in E_1} T_e
  \le o(t^*).
\]
Điều này mâu thuẫn với $o(t^*)=0$. Vậy $C^*$ là chu trình tối ưu.

Nếu $t^*$ là đáp án tối ưu, ta có tính chất
\[
\begin{cases}
  o(t)=0 \Longleftrightarrow t=t^*,\\
  o(t)<0 \Longleftrightarrow t>t^*,\\
  o(t)>0 \Longleftrightarrow t<t^*.
\end{cases}
\tag{1}
\]
Từ đây thuật toán xuất hiện tự nhiên. Bắt đầu với một khoảng chứa $t^*$, chẳng
hạn
\[
  t_l=\min_{e\in E}\frac{C_e}{T_e},\qquad
  t_h=\max_{e\in E}\frac{C_e}{T_e}.
\]
Mỗi lần lấy trung điểm $t$ của $t_l$ và $t_h$, rồi tính $o(t)$.

Cách tính $o(t)$ như sau: với mỗi cạnh $e\in E$, đặt trọng số mới
\[
  W_e=C_e-tT_e.
\]
Dùng Floyd để tìm chu trình có tổng trọng số lớn nhất trong đồ thị có hướng.
Giá trị lớn nhất ấy chính là $o(t)$.

Theo tính chất (1), ta cập nhật $t_l$ hoặc $t_h$, tiếp tục nhị phân cho đến
khi đạt độ chính xác yêu cầu. Nếu số lần nhị phân là $K$, độ phức tạp thời
gian là $\Oh(Kn^3)$. Tuy đây không phải một thuật toán đa thức nghiêm ngặt,
với giới hạn của đề bài nó vẫn tìm được đáp án rất nhanh.

\section{Ví dụ 3: Building Towers}

\subsection{Tóm tắt đề bài}

Có $N$ khối lập phương, cần dùng chúng để xây tháp. Mỗi tháp có $H$ tầng; tầng
dưới cùng chứa $M$ khối. Với mỗi tầng phía trên, số khối ở tầng đó phải đúng
bằng số khối ở tầng ngay dưới cộng $1$ hoặc trừ $1$.

Trong tài liệu gốc có một hình minh họa cho một tháp hợp lệ với
$H=6$, $M=2$, $N=13$. Nhiệm vụ là tính số tháp khác nhau có thể xây được; lưu
ý không bắt buộc dùng hết mọi khối. Hai tháp được xem là khác nhau nếu tồn
tại một tầng $i$ mà số khối của tầng thứ $i$ trong hai tháp khác nhau. Mỗi
tháp được biểu diễn bằng một dãy $H$ số nguyên dương từ đáy lên đỉnh. Ngoài
ra cần trả lời các truy vấn: in tháp có thứ tự từ điển thứ $K$.

Dữ liệu vào gồm $N,H,M$ với $N\le 32767$, $H\le 60$, $M\le 10$, sau đó là các
truy vấn $K$; dòng cuối cùng chứa $-1$. Dữ liệu ra gồm tổng số tháp khác nhau
và, với mỗi truy vấn, dãy biểu diễn tháp tương ứng.

\subsection{Phân tích ban đầu}

Khi vừa nhìn đề, theo quán tính ta rất dễ nghĩ đến quy hoạch động cổ điển.
Gọi $F[i,j,k]$ là số phương án khi tầng hiện tại có $i$ khối, còn $j$ tầng và
còn có thể dùng $k$ khối. Khi đó có thể viết công thức
\[
  F[i,j,k]=F[i+1,j+1,k-i]+F[i-1,j+1,k-i].
\]
Điều kiện biên là $F[M,H,N]=1$, các trạng thái còn lại ban đầu bằng $0$.

Dễ thấy quy hoạch động này liên quan tới khoảng $N\cdot M\cdot K$ trạng thái,
trường hợp lớn nhất vào khoảng $60\cdot 70\cdot 2400\approx 10^7$. Nếu mỗi
trạng thái dùng một số \texttt{double}, bộ nhớ cần ít nhất khoảng $80$ MB.
Hơn nữa, do cần in tháp thứ $K$ theo thứ tự từ điển, ta không thể dùng tối ưu
cuốn chiếu. Cả thời gian lẫn bộ nhớ đều khó chấp nhận, nên hướng quy hoạch
động thuận bị bế tắc.

\subsection{Dùng tư duy nghịch hướng}

Quy hoạch động từ dưới lên gây áp lực lớn về thời gian và không gian. Ta thử
``đảo chiều'' xử lý và xem xét tìm kiếm từ trên xuống.

Tìm kiếm thường bị xem là đồng nghĩa với kém hiệu quả và hàm mũ. Nhưng nếu
thêm ghi nhớ hóa, nó lại trở thành một quá trình quy hoạch động nghịch hướng.
Đảo thứ tự quy hoạch có hai lợi ích.

Thứ nhất, nhìn bài toán từ trên xuống giúp ta có cái nhìn tổng thể hơn, thuận
lợi cho việc quyết định tìm kiếm thế nào, tìm gì trước, tìm gì sau và cắt
tỉa ra sao. Trong bài này, tìm kiếm sẽ không đi vào nhiều trạng thái vô lý,
chẳng hạn số khối rõ ràng không đủ hoặc sai khác về tính chẵn lẻ. Với các bài
tối ưu, tìm kiếm có ghi nhớ cũng phối hợp tốt với cắt tỉa theo tính tối ưu,
tránh nhiều trạng thái không có giá trị, điều mà quy hoạch động thuận khó làm.

Thứ hai, trong các bài đếm tương tự bài này, có thể dùng ghi nhớ hóa một phần:
chỉ lưu những trạng thái dễ bị tìm lại nhiều lần. Như vậy số trạng thái phải
lưu giảm xuống và tốc độ tra cứu tăng lên.

Khi cài đặt, dùng bộ ba $(N,H,M)$ để biểu diễn một trạng thái: tầng hiện tại
có $M$ khối, còn $H$ tầng, và còn có thể dùng $N$ khối. Ký hiệu
$F(N,H,M)$ là số phương án của trạng thái ấy. Ta lưu kết quả bằng bảng băm và
thêm một số tối ưu sau.

\begin{itemize}
  \item Tính trước số khối ít nhất và nhiều nhất mà trạng thái hiện tại còn
  cần. Nếu $N$ nhỏ hơn nhu cầu tối thiểu thì đáp án rõ ràng là $0$; nếu $N$
  lớn hơn nhu cầu tối đa, đặt lại $N$ bằng nhu cầu tối đa để tránh lưu trùng
  các trạng thái tương đương.
  \item Nếu $N$ không nhỏ hơn nhu cầu tối đa và $H\le M$, thì mọi lựa chọn
  lên/xuống ở $H-1$ bước còn lại đều hợp lệ, nên đáp án là $2^{H-1}$.
  \item Để giảm lượng trạng thái lưu trữ, đặt một ngưỡng lưu: chỉ khi
  $F(N,H,M)$ không vượt quá ngưỡng ấy mới đưa trạng thái vào bảng băm.
\end{itemize}

Vì sao tối ưu cuối cùng có hiệu quả? Cây trạng thái của tìm kiếm mở rộng theo
dạng hàm mũ. Càng xuống sâu, số trạng thái càng nhiều và tần suất tra cứu
càng cao; gần gốc, trạng thái phân bố thưa hơn, và khi $F(N,H,M)$ lớn thì tần
suất truy cập chúng cũng giảm. Do đó ta có thể chấp nhận tìm lại các trạng
thái lớn ấy để giảm bộ nhớ và tăng tốc tra cứu bảng băm.

Hiệu quả thực nghiệm của ``quy hoạch động nghịch hướng'' như sau.
\[
\begin{array}{c|c|c|r|r|r}
N & H & M & \text{Số trạng thái của QHĐ thuận}
  & \text{Số trạng thái lưu trong bảng băm}
  & \text{Tỉ lệ}\\ \hline
1000 & 40 & 10 & 1200000 & 3055  & 392.80\\
1500 & 50 & 10 & 2625000 & 5371  & 488.74\\
1200 & 60 & 10 & 2880000 & 49875 & 57.74\\
1500 & 60 & 10 & 3600000 & 38018 & 94.69
\end{array}
\]
Ngay trong trường hợp xấu nhất, quy hoạch động nghịch hướng cũng chỉ cần xử
lý khoảng một phần năm mươi số trạng thái. Đến đây, bài toán đã được giải
trọn vẹn.

\section{Tổng kết}

Ví dụ thứ nhất là tư tưởng chuyển sang phần bù: khi không biết xử lý tập ban
đầu thế nào, ta đếm phần bù để nhìn bài toán từ mặt ngược lại. Ví dụ thứ hai
không nhìn mặt ngược của bài toán, mà nhìn mặt ngược của một sự thật: hiện ta
chưa biết đáp án, vậy nếu đã biết đáp án thì điều gì sẽ xảy ra? Ví dụ thứ ba
đảo chiều thứ tự quy hoạch khi quy hoạch động cổ điển không giải nổi, từ đó
thu được quy hoạch động nghịch hướng, tránh nhiều trạng thái vô hiệu và cải
thiện mạnh cả thời gian lẫn không gian.

Ba ví dụ trên chỉ là một vài ứng dụng đơn giản của tư duy nghịch hướng. Tinh
thần ``khó thuận thì xét ngược'' giúp ta phá vỡ quán tính tư duy, lấy lùi làm
tiến, tránh chỗ mạnh và đánh vào chỗ yếu của bài toán. Khi con đường quen
thuộc không còn lối, đổi hướng suy nghĩ có thể đem lại lời giải rất tự nhiên.

\section*{Lời cảm ơn}

Tác giả đặc biệt cảm ơn bạn Zhou Yuan từ Trường Trung học số 1 Vu Hồ, An Huy,
cũng như tất cả những người đã giúp đỡ tác giả.

\section*{Tài liệu tham khảo}

\begin{enumerate}
  \item Liu Rujia và Huang Liang, \textit{Nghệ thuật thuật toán và thi đấu tin
  học}, Nhà xuất bản Đại học Thanh Hoa.
  \item Gunnar W. Klau và các đồng tác giả, \textit{The Fractional
  Prize-Collecting Steiner Tree Problem On Trees}.
  \item Ren Chunong, \textit{Bàn về tư duy nghịch hướng}.
\end{enumerate}

\section*{Phụ lục về chương trình}

Mã nguồn cho ví dụ 3, bài \textit{Building Tower}, nằm trong tệp
\texttt{Tower.cpp} đi kèm tài liệu gốc.

\appendix
\section{Nguyên đề của các ví dụ}

\subsection{Example 1: Dinner Is Ready}

Dinner is ready! WishingBone's family rush to the table. WishingBone's mother
has prepared many delicious bones. WishingBone has several brothers, namely
WishingTwoBones, WishingThreeBones. WishingBone wants at least one bone,
WishingTwoBones two and WishingThreeBones three. :-P But as they are not too
greedy (or they want to keep shape), they decide that they want at most twice
of their minimal requirement. Now let's suppose mother has prepared 7 bones,
one way to distribute them is 2-2-3, the other two ways are 1-3-3 and 1-2-4.
Of course mother doesn't want to waste any bones, so if she had prepared 13
bones, she would end up without any feasible distribution.

As curious as WishingBone, mother wants to know how many different ways she can
distribute those bones.

\paragraph{Input.}
The first line of input is a positive integer $N \le 100$, which is the number
of test cases. The first line of each case contains two integers $n$ and $m$
($0<n\le 8$, $0<m$), where $n$ is the number of children and $m$ is the number
of bones mother has prepared. The next $n$ lines have two integers
\texttt{mini} and \texttt{maxi} ($0\le\texttt{mini}\le\texttt{maxi}\le m$),
which are the minimal and maximal requirement of the $i$-th child.

\paragraph{Output.}
$N$ integers which are the numbers of ways to distribute the $m$ bones, one
per line.

\paragraph{Sample Input.}
\begin{verbatim}
2
3 7
1 2
2 4
3 6
3 13
1 2
2 4
3 6
\end{verbatim}

\paragraph{Sample Output.}
\begin{verbatim}
3
0
\end{verbatim}

\subsection{Example 2: Greedy Path}

There are $N$ towns and $M$ routes between them. Recently, some new travel
agency was founded for servicing tourists in these cities. We know cost which
tourist has to pay, when traveling from town $i$ to town $j$ which equals to
$C_{ij}$ and time needed for this travel, $T_{ij}$. There are many tourists
who want to use this agency because of very low rates for travels, but agency
still has only one bus. Head of agency decided to organize one big travel
route to gain maximal possible amount of money. Scientists of the company
offer to find such a cyclic path $G$, when greedy function $f(G)$ will be
maximum. Greedy function for some path is calculated as total cost of the path
divided by total time of path. But nobody can find this path, and Head of the
company asks you to help him in solving this problem.

\paragraph{Input.}
There are two integers $N$ and $M$ on the first line of input file
($3\le N\le 50$). Next $M$ lines contain routes information, one route per
line. Every route description has format $A,B,C_{ab},T_{ab}$, where $A$ is
starting town for route, $B$ is ending town for route, $C_{ab}$ is cost of
route and $T_{ab}$ is time of route ($1\le C_{ab}\le 100$,
$1\le T_{ab}\le 100$, $A\ne B$). Note that order of towns in route is
significant: route $(i,j)$ is not equal to route $(j,i)$. There is at most one
route in one direction between any two towns.

\paragraph{Output.}
You must output requested path $G$ in the following format. On the first line
of output file you must output $K$, the number of towns in the path
($2\le K\le 50$); on the second line, numbers of these towns in order of
passing them. If there are many such ways, output any one of them; if there
are no such ways, output \texttt{"0"}.

\paragraph{Sample Input.}
\begin{verbatim}
4 5
1 2 5 1
2 3 3 5
3 4 1 1
4 1 5 2
2 4 1 10
\end{verbatim}

\paragraph{Sample Output.}
\begin{verbatim}
4
1 2 3 4
\end{verbatim}

\subsection{Example 3: Building Towers}

Time Limit: 2 second. Memory Limit: 1000 KB.

There are $N$ cubes in a toy box which has 1-unit height; the width is double
the height. The teacher organizes a tower-building game. The tower is built by
the cubes. The height of the tower is $H$ levels. The bottom of the tower
contains $M$ cubes; and for all above level, each must contain a number of
cubes which is exactly 1 less than or greater than the number of cubes of the
level right below it.

The original document shows an example of a tower with $H=6$, $M=2$, $N=13$.

Your task is to determine how many different towers can be there. Two towers
are considered different if there is at least one number $i$ ($1<i\le H$) so
that the $i$-th level of one tower contains a different number of cubes to the
$i$-th level of the other tower. Each tower is represented by an array of $H$
positive integers which determines the structure of the tower from the bottom
level up to the top level. Those arrays are sorted ascendingly.

\paragraph{Input.}
The first line contains three positive numbers $N,H,M$ ($N\le 32767$,
$H\le 60$, $M\le 10$). Each following line, except the last, contains an
integer $K$. The last line contains number $-1$.

\paragraph{Output.}
The first line is the total of different towers that can be founded. Each
following line contains an array that represents the structure of the tower
corresponding to the number $K$ from input.

\paragraph{Sample Input.}
\begin{verbatim}
22 5 4
1
10
-1
\end{verbatim}

\paragraph{Sample Output.}
\begin{verbatim}
10
4 3 2 1 2
4 5 4 5 4
\end{verbatim}

\paragraph{Hint.}
Numbers which are on the same line should be separated by at least one space.
You do not have to use all $N$ cubes.

This problem also appears as Zhejiang University Online Judge 1442, Saratov
State University Online Judge 236, and Ural Online Judge 1148. One related
note in the source says that this is a classical combinatorics problem that
can be solved by the divider principle.

\end{document}
